\documentclass[11pt,a4paper]{article}
\usepackage[utf8]{inputenc}
\usepackage[a4paper,top=2.3cm,bottom=2.6cm,left=2.5cm,right=2.5cm,headsep=0.6cm,footskip=1.1cm]{geometry}
\usepackage{amsmath,amssymb}
\usepackage{enumitem}
\usepackage{fancyhdr}
\usepackage{titlesec}
\titlespacing*{\section}{0pt}{1.3em}{0.2em}
\usepackage{hyperref}
\hypersetup{colorlinks=true,linkcolor=black,urlcolor=black,
  pdftitle={FZT 203E Problem Set 3},pdfauthor={Mehmet Özkan}}

\newcommand{\Z}{\mathbb{Z}}
\newcommand{\R}{\mathbb{R}}
\newcommand{\C}{\mathbb{C}}
\newcommand{\gen}[1]{\langle #1\rangle}
\DeclareMathOperator{\sgn}{sgn}
\DeclareMathOperator{\tr}{tr}
\DeclareMathOperator{\diag}{diag}
\DeclareMathOperator{\spn}{span}
\newcommand{\ket}[1]{\lvert #1\rangle}
\newcommand{\up}{{\uparrow}}
\newcommand{\dn}{{\downarrow}}
\newcommand{\MM}{\mathsf{M}}

\pagestyle{fancy}
\fancyhf{}
\renewcommand{\headrulewidth}{0pt}
\fancyfoot[L]{\small FZT 203E $\cdot$ Problem Set 3}
\fancyfoot[R]{\small\thepage}

\setlength{\parindent}{0pt}
\setlength{\parskip}{0.55em}

\newcommand{\problem}[1]{\par\vspace{0.9em}\noindent{\bfseries #1}\par\vspace{0.1em}}
\setlist[enumerate,1]{label=\textbf{(\alph*)},leftmargin=2.2em,itemsep=0.35em,topsep=0.2em}
\setlist[itemize]{label=$\bullet$,itemsep=0.25em,topsep=0.2em}

\begin{document}

{\small FZT 203E $\cdot$ Theory of Groups in Physics $\cdot$ Fall 2026 \hfill Istanbul Technical University\\
Instructor: Prof.\ Mehmet Özkan}

\vspace{0.9em}
{\LARGE\bfseries Problem Set 3}

\vspace{0.5em}
Problem Set 3 covers Lectures 7--10: linear algebra, representations, Schur's lemma, the Great Orthogonality Theorem and characters. Everything from Problem Sets 1 and 2 may be used.

\textbf{Due:} Wednesday, 25 November 2026, at the start of the lecture.\\
\textbf{Reading:} lecture notes for Weeks 4 and 5; Zee, chapters II.1--II.3. Zee III.1 (Bloch's theorem) goes with Problem 2, and Zee III.2 (harmonic motion and zero modes) with Problem 1.

Justify every answer; a correct answer without reasoning earns partial or no credit.

\section*{\large Conventions}
These are the conventions of the lectures and of Problem Sets 1 and 2.
\begin{itemize}
  \item Products are read right to left: in $gh$, the element $h$ acts first. A permutation $p$ sends position $i$ to position $p(i)$, and $(12)(23)$ means first $(23)$, then $(12)$. For the triangle, $r = (123)$, $s = (23)$ and $rs = (12)$.
  \item Vectors are columns. A representation is a homomorphism $D\colon G \to GL(d, \C)$, $D(gh) = D(g)D(h)$; representations are finite dimensional over $\C$ unless a real space is specified.
  \item For a finite group, $N = |G|$; the conjugacy classes are $C_1, \dots, C_k$, with $n_i = |C_i|$ elements. The character of $D$ is $\chi_D(g) = \tr D(g)$. For functions on $G$,
  \[
    \langle f, h\rangle_G = \frac{1}{N}\sum_{g \in G} f(g)^* h(g) = \frac{1}{N}\sum_{i=1}^{k} n_i\, f(C_i)^* h(C_i),
  \]
  the second form for class functions. The multiplicity of the irreducible representation $D^{(\alpha)}$ in $D$ is $m_\alpha = \langle \chi_\alpha, \chi_D\rangle_G$, and $D$ is irreducible exactly when $\langle \chi_D, \chi_D\rangle_G = 1$.
  \item The plane representation of $D_3 \cong S_3$, with vertex 1 at $(1, 0)$, is
  \[
    D_E(r) = R = \begin{pmatrix} -\tfrac12 & -\tfrac{\sqrt3}{2} \\[2pt] \tfrac{\sqrt3}{2} & -\tfrac12 \end{pmatrix}, \qquad
    D_E(s) = F = \begin{pmatrix} 1 & 0 \\ 0 & -1 \end{pmatrix}.
  \]
  The classes of $S_3$ are $\{e\}$, the two 3-cycles $\{r, r^2\}$ (rotations) and the three transpositions $\{s, rs, r^2 s\}$ (reflections). Its irreducible representations are the trivial representation $1$, the sign representation $\sgn$ and $E$, with the characters
  \[
    \begin{array}{c|ccc}
      & \{e\} & \{r, r^2\} & \{s, rs, r^2 s\} \\ \hline
      1 & 1 & 1 & 1 \\
      \sgn & 1 & 1 & -1 \\
      E & 2 & -1 & 0
    \end{array}
  \]
  \item You may use the results of the lectures, among them the finite-group unitarity theorem, both forms of Schur's lemma, the Great Orthogonality Theorem, the orthonormality of irreducible characters, $\sum_\alpha d_\alpha^2 = N$, and that the number of inequivalent irreducible representations equals the number of conjugacy classes.
\end{itemize}

\section*{\large Problems}

\problem{1. The vibrations of carbon dioxide}
The carbon dioxide molecule is straight, O=C=O, with the carbon atom (mass $M$) between two oxygen atoms (mass $m$). Consider motion along the axis only: $x_1$, $x_2$ and $x_3$ are the displacements of the three atoms from equilibrium, and each bond acts as a spring with constant $k$:
\[
  m\ddot x_1 = -k(x_1 - x_2), \qquad M\ddot x_2 = -k(x_2 - x_1) - k(x_2 - x_3), \qquad m\ddot x_3 = -k(x_3 - x_2).
\]
\begin{enumerate}
  \item Write these equations as $\MM\ddot x = -Kx$, with $x = (x_1, x_2, x_3)^T$ and the mass matrix $\MM = \diag(m, M, m)$, and find the symmetric matrix $K$. Show that $x(t) = a\cos\omega t$, with a fixed vector $a \ne 0$, solves the equations exactly when $Ka = \omega^2\MM a$, that is, when $a$ is an eigenvector of $A = \MM^{-1}K$ with eigenvalue $\omega^2$. Show that $a^T K a = k\bigl[(a_1 - a_2)^2 + (a_2 - a_3)^2\bigr]$, and deduce that $\omega^2 \ge 0$ for every real mode $a$.
  \item $A$ is not symmetric when $M \ne m$. Show that $\langle a, b\rangle_\MM = a^T \MM b$ is an inner product on $\R^3$ for which $\langle a, Ab\rangle_\MM = \langle Aa, b\rangle_\MM$, and deduce that two modes with different frequencies are orthogonal for $\langle\cdot, \cdot\rangle_\MM$.
  \item The inversion through the carbon atom exchanges the two oxygen atoms and reverses the direction of the axis, so it maps a displacement $x$ to $Px = (-x_3, -x_2, -x_1)^T$. Write $P$ as a matrix, and show that $P^2 = I$ and that $P$ commutes with $\MM$ and $K$, and hence with $A$. Thus $\{I, P\}$ is a representation of $\Z_2$ on $\R^3$. Use its character to find how many times the trivial and the sign representations of $\Z_2$ appear in it. Find the subspaces of $\R^3$ on which $P = +1$ and $P = -1$, and show that $A$ maps each of them into itself.
  \item Find the three frequencies and modes: on the line where $P = +1$ directly, and on the plane where $P = -1$ by writing $A$ there as a $2 \times 2$ matrix in the basis $(1, 0, 1)^T$, $(0, 1, 0)^T$. Describe each mode in words. Why is one frequency zero? Check that the other two modes are orthogonal to it for $\langle\cdot, \cdot\rangle_\MM$, and explain what this means for the total momentum of the molecule.
  \item The oxygen atoms carry a small charge $-q$ each and the carbon atom $+2q$, so a displacement gives the molecule the electric dipole moment $d(x) = q(2x_2 - x_1 - x_3)$ along its axis. Show that $d(Px) = -d(x)$, and deduce that $d = 0$ at all times in a mode with $P = +1$. Only a vibration that makes the dipole oscillate absorbs infrared light: which of your modes does? This asymmetric stretch of CO$_2$ absorbs near a wavelength of $4.3\,\mu$m, and its bending vibrations (sideways, not in this model) near $15\,\mu$m, close to the peak of the infrared radiation emitted by the Earth. Why can a molecule of two identical atoms, like N$_2$ or O$_2$, the main gases of the air, not absorb infrared light in this way?
\end{enumerate}

\problem{2. Why polyacetylene is not a metal}
Polyacetylene is a long chain of carbon atoms, each bonded to one hydrogen atom, with alternating double and single bonds between the carbons. In the simplest (H\"uckel) model of its $\pi$ electrons, take a ring of $N \ge 3$ cells, each containing two carbon atoms $A_n$ and $B_n$, with $n = 1, \dots, N$ counted modulo $N$ (so $A_{N+1} = A_1$). A state of an electron is a vector in $\C^{2N}$ with basis vectors $e_{A_n}$, $e_{B_n}$, and its energies are the eigenvalues of the Hamiltonian
\[
  He_{A_n} = -t_1 e_{B_n} - t_2 e_{B_{n-1}}, \qquad He_{B_n} = -t_1 e_{A_n} - t_2 e_{A_{n+1}},
\]
where $t_1 > 0$ is the hopping amplitude across the double bond $A_nB_n$ and $t_2 > 0$ across the single bond $B_nA_{n+1}$. The translation by one cell is $Te_{A_n} = e_{A_{n+1}}$, $Te_{B_n} = e_{B_{n+1}}$.
\begin{enumerate}
  \item Show that $H$ is real and symmetric, that $T^N = I$, and that $HT = TH$. Thus the powers of $T$ form a representation of $\Z_N$ on $\C^{2N}$, and $H$ commutes with every matrix in it.
  \item Use Schur's lemma to show that every irreducible representation of an Abelian group is one-dimensional. Show that $\Z_N = \{e, a, \dots, a^{N-1}\}$ has exactly $N$ inequivalent irreducible representations, $D(a^j) = \lambda^j$ with $\lambda^N = 1$.
  \item For $k = 2\pi l/N$ with $l = 0, 1, \dots, N - 1$, let
  \[
    u^A_k = \sum_{n=1}^{N} e^{ikn} e_{A_n}, \qquad u^B_k = \sum_{n=1}^{N} e^{ikn} e_{B_n}.
  \]
  Show that $Tu^A_k = e^{-ik}u^A_k$ and $Tu^B_k = e^{-ik}u^B_k$, and that the $2N$ vectors $u^A_k$, $u^B_k$ are orthogonal. (Hint: sum a geometric series.) Conclude that $\C^{2N}$ is the direct sum of the planes $V_k = \spn(u^A_k, u^B_k)$, and that each one-dimensional representation of $\Z_N$ appears in it exactly twice.
  \item Show that $H$ maps each $V_k$ into itself, with the matrix
  \[
    h(k) = -\begin{pmatrix} 0 & t_1 + t_2e^{-ik} \\ t_1 + t_2e^{ik} & 0 \end{pmatrix}
  \]
  in the basis $(u^A_k, u^B_k)$, and find its eigenvalues $E_\pm(k)$. The $2N \times 2N$ problem has become $N$ problems of size $2 \times 2$: this is Bloch's theorem. As a check, take $N = 3$ and $t_1 = t_2 = t$, a ring of six equal bonds like benzene, and list the six energies with their degeneracies.
  \item In the limit $N \to \infty$, which we take from now on, the allowed $k$ fill the interval $(-\pi, \pi]$ densely, and the energies form two bands. Find the range of each band. Every carbon atom gives one $\pi$ electron, and each level holds two electrons (spin up and spin down): show that the $2N$ electrons exactly fill the lower band. Show that for $t_1 = t_2$ the two bands touch, so an arbitrarily small energy lifts an electron into an empty level and the chain conducts like a metal, while for $t_1 \ne t_2$ there is an energy gap $2|t_1 - t_2|$ between them. Real polyacetylene has alternating bonds and starts to absorb light at about $1.5$\,eV. Taking this as the gap, what is $|t_1 - t_2|$? (The discovery that polyacetylene conducts well when doped with iodine was recognized by the 2000 Nobel Prize in Chemistry.)
\end{enumerate}

\problem{3. Degenerate modes of a square drum}
The skin of a square drum covers $-1 \le x, y \le 1$ and is fixed along the edges. Its normal modes are vibrations $u(x, y)\cos\omega t$ with
\[
  u_{nm}(x, y) = f_n(x)f_m(y), \qquad f_n(x) = \sin\frac{n\pi(x + 1)}{2}, \qquad n, m = 1, 2, 3, \dots,
\]
and frequencies $\omega_{nm} = \frac{\pi c}{2}\sqrt{n^2 + m^2}$, where $c$ is the speed of waves on the skin. (A quantum particle in a square box has the same wavefunctions, with energies proportional to $n^2 + m^2$.) The symmetry group $D_4$ of the square acts on the plane by the quarter-turn and the reflection in the $x$-axis,
\[
  D(r) = \begin{pmatrix} 0 & -1 \\ 1 & 0 \end{pmatrix}, \qquad D(s) = \begin{pmatrix} 1 & 0 \\ 0 & -1 \end{pmatrix},
\]
and on functions by $(O_g u)(x, y) = u\bigl(D(g)^{-1}(x, y)^T\bigr)$; you may use that $O_{gh} = O_g O_h$. The character table of $D_4$ is
\[
  \begin{array}{c|ccccc}
    & \{e\} & \{r^2\} & \{r, r^3\} & \{s, r^2 s\} & \{rs, r^3 s\} \\ \hline
    1_{++} & 1 & 1 & 1 & 1 & 1 \\
    1_{+-} & 1 & 1 & 1 & -1 & -1 \\
    1_{-+} & 1 & 1 & -1 & 1 & -1 \\
    1_{--} & 1 & 1 & -1 & -1 & 1 \\
    E & 2 & -2 & 0 & 0 & 0
  \end{array}
\]
where $1_{\epsilon\epsilon'}$ is the one-dimensional representation with $r \mapsto \epsilon$ and $s \mapsto \epsilon'$, and $E$ is the representation $D$ itself.
\begin{enumerate}
  \item Show that $f_n(-x) = (-1)^{n+1}f_n(x)$. Show that
  \[
    O_r u_{nm} = (-1)^{m+1}u_{mn}, \qquad O_s u_{nm} = (-1)^{m+1}u_{nm},
  \]
  and find $O_{r^2}u_{nm}$ and $O_{rs}u_{nm}$. Deduce that every $O_g$ maps the modes of each frequency to modes of the same frequency, so that these modes span a representation of $D_4$.
  \item For $n \ne m$, the modes $u_{nm}$ and $u_{mn}$ have the same frequency. Find the character of the representation of $D_4$ on $W_{nm} = \spn(u_{nm}, u_{mn})$ on the five classes. Show that $W_{nm}$ is irreducible, and equivalent to $E$, exactly when $n + m$ is odd.
  \item Let $n \ne m$ with $n + m$ even. Decompose $W_{nm}$ with characters, find its two invariant lines, and say which one-dimensional representation acts on each. Treat the cases of $n$ and $m$ both odd and both even separately.
  \item A small change of the drum that keeps all the symmetries of the square, for example a small weight glued to its centre, shifts the frequencies. For the modes of one frequency, the first-order shifts of their frequencies are the eigenvalues of a real symmetric matrix $V$ that acts on their span and commutes with every $O_g$, and its eigenvectors are the new modes. Use Schur's lemma to show that $u_{12}$ and $u_{21}$ keep a common frequency, while $u_{13}$ and $u_{31}$ may split, and that the new modes are then $u_{13} \pm u_{31}$. Check this for the weight at the centre: a mode with a node at the centre is not affected, and a mode that moves there is slowed down. (Use $\sin\frac{\pi}{2} = 1$, $\sin\pi = 0$ and $\sin\frac{3\pi}{2} = -1$.) The degeneracy of $u_{13}$ and $u_{31}$ is called accidental: the symmetry of the square does not require it.
  \item The smallest value of $n^2 + m^2$ shared by three modes is $50 = 1^2 + 7^2 = 5^2 + 5^2$. Without computing, explain why the symmetry of the square cannot be the reason for this threefold degeneracy. Then find the character of the representation on the span of $u_{17}$, $u_{71}$, $u_{55}$, and decompose it.
\end{enumerate}

\problem{4. Three spins on a triangle}
Three spin-one-half particles, for example the unpaired electrons of three magnetic ions at the corners of a triangle, sit at the sites $1, 2, 3$. Each spin is up ($\up$) or down ($\dn$), so a state is a vector in $\C^8$ with the basis $\ket{s_1s_2s_3}$, $s_i \in \{\up, \dn\}$. A permutation $p \in S_3$ moves the spin at site $i$ to site $p(i)$:
\[
  P(p)\ket{s_1s_2s_3} = \ket{s'_1s'_2s'_3}, \qquad s'_{p(i)} = s_i.
\]
For example, $P\bigl((12)\bigr)\ket{\up\dn\dn} = \ket{\dn\up\dn}$.
\begin{enumerate}
  \item Show that $P(pq) = P(p)P(q)$, and that $\chi_P(p)$ is the number of basis states left unchanged by $P(p)$. Find $\chi_P$ on the three classes of $S_3$, and decompose $P$ into irreducible representations.
  \item Permutations do not change the number of up spins, so the subspaces with 3, 2, 1 and 0 up spins (of dimensions 1, 3, 3 and 1) are invariant. Decompose each of them, and check that the results add up to (a). In the subspace with two up spins, find the line on which $S_3$ acts trivially and an invariant complement on which it acts by $E$.
  \item Three electrons in the same orbital, such as the 1s orbital of an atom, have a wavefunction that is symmetric in their positions, so the Pauli principle requires their spin state $\psi$ to be antisymmetric: $P(p)\psi = \sgn(p)\psi$ for every $p \in S_3$. Show that no such $\psi \ne 0$ exists. The third electron of a lithium atom therefore cannot occupy the 1s orbital; the lowest orbital still available to it is 2s.
  \item Let $K = P\bigl((12)\bigr) + P\bigl((23)\bigr) + P\bigl((13)\bigr)$. Show that $K$ commutes with every $P(p)$. More generally, for any representation $D$ of $S_3$ show that $K_D = D\bigl((12)\bigr) + D\bigl((23)\bigr) + D\bigl((13)\bigr)$ commutes with every $D(p)$, and use Schur's lemma to show that for irreducible $D$ of dimension $d$, $K_D = \lambda I$ with $\lambda = 3\chi_D(\tau)/d$, where $\tau$ is a transposition. Find $\lambda$ for $1$, $\sgn$ and $E$.
  \item The Heisenberg model of the three ions has the Hamiltonian $H = J(\vec S_1\cdot\vec S_2 + \vec S_2\cdot\vec S_3 + \vec S_3\cdot\vec S_1)$, where $\vec S_i$ is the spin of the particle at site $i$. For spin one-half, Dirac's identity $\vec S_i\cdot\vec S_j = \tfrac12 P\bigl((ij)\bigr) - \tfrac14$ (with $\hbar = 1$) turns it into
  \[
    H = \frac{J}{2}K - \frac{3J}{4},
  \]
  and you may take this as the definition of $H$. Without diagonalizing an $8 \times 8$ matrix, find the energies and their degeneracies. For $J > 0$ every pair of neighbours prefers opposite spins, which cannot happen for all three pairs of a triangle: the magnet is frustrated. Show that its lowest level is fourfold degenerate. A magnetic field along the $z$-axis adds $-hS_z$ to $H$, where $S_z$ is half the number of up spins minus half the number of down spins. Show that for $0 < h < \tfrac32 J$ the lowest level is still twofold degenerate, and explain with Schur's lemma why the field cannot remove this degeneracy.
\end{enumerate}

\end{document}
